% 03_phuongphap.tex - Phần 3: Phương pháp và giải pháp
\section{Phương pháp và giải pháp: Tam giác liên ngành}

Để đối phó với sự phức tạp của các bộ máy gợi ý tỷ đô, kiến trúc Algorithmic Circuit Breaker v3.0 áp dụng tam giác ứng phó liên ngành (The interdisciplinary triumvirate), quy đổi các sự kiện tâm lý thành lý thuyết điều khiển, xử lý ngôn ngữ tự nhiên và học tăng cường.

\subsection{Cơ chế đối phó thiết kế gây nghiện: Mô hình Rescorla-Wagner}

Các thuật toán mạng xã hội hiện đại cực đại hóa sự gắn kết (engagement) bằng cách thao túng hệ thống dopamine thông qua các lịch trình phần thưởng biến thiên. Để kiến trúc hóa biện pháp đối phó, hệ thống lượng hóa sự kỳ vọng nội tâm bằng cách mô hình hóa người dùng như một tác tử học tăng cường (Reinforcement Learning Agent).

Cơ sở sinh học thần kinh cho cách tiếp cận này bắt nguồn từ khám phá nền tảng của Schultz và cộng sự (1997) \cite{schultz1997}, trong đó chứng minh rằng các nơ-ron tiết dopamine không mã hóa giá trị phần thưởng tuyệt đối, mà chỉ phản ứng với sai lệch dự đoán phần thưởng (Reward Prediction Error - RPE). Tuy nhiên, nguyên lý sinh học vi mô này chỉ thiết lập nền tảng cho phương trình RPE, hoàn toàn không đủ cơ sở để trực tiếp ngoại suy hay định lượng hành vi tiêu thụ nội dung số vĩ mô.

Để chứng minh tính hợp lệ của việc dùng mô hình Học tăng cường để đo lường mức độ nghiện của con người trên không gian mạng, hệ thống được thiết kế hoàn toàn dựa trên khuôn khổ của Tâm thần học Tính toán (Computational Psychiatry). Cụ thể, nghiên cứu của Kato và cộng sự (2023) về các mô hình tính toán cho chứng nghiện hành vi \cite{kato2023} đã chính thức công nhận Học tăng cường là khuôn khổ chuẩn mực trong y văn để lượng hóa các rối loạn bốc đồng, điển hình như Rối loạn chơi game trên Internet (Internet Gaming Disorder). Việc biểu diễn tín hiệu RPE cho phép chúng tôi cung cấp một thước đo định lượng đại diện cho trạng thái khao khát (craving) sinh ra từ sự thao túng phần thưởng.

Trong khối kiến trúc này, mức độ nghiện được cập nhật qua từng bước thời gian dưới dạng RPE theo phương trình Rescorla-Wagner:

\begin{equation}
    V(t+1) = V(t) + \alpha [R(t) - V(t)]
\end{equation}

Với $RPE = R(t) - V(t)$ đóng vai trò là biến trạng thái phản hồi. Do tính chất trừu tượng của giao diện, hệ thống trích xuất vận tốc cuộn (scroll velocity) làm biến đại diện hành vi (behavioral proxy) cho tín hiệu tìm kiếm phần thưởng $R(t)$. Tốc độ cuộn càng cao tỷ lệ thuận với cường độ khao khát bị kích hoạt ở Hệ thống vô thức. Để duy trì tính liên tục của động học thần kinh, một cơ chế tích phân rò rỉ (leaky integrator) được áp dụng: khi $RPE \le 0$, mức độ khao khát không triệt tiêu theo hàm bước (step function) mà suy thoái từ từ theo hàm mũ với hệ số $\gamma = 0.95$. Giá trị RPE tích lũy này sau đó được cấp làm tín hiệu nhiễu cho khối điều khiển.

\subsection{Cơ chế đối phó thao túng dư luận: Shift-Left Local NLP Heuristics}

\subsubsection{Xử lý Ngôn ngữ Tự nhiên Cục bộ (Shift-Left Local NLP Heuristics)}

Nhằm tuân thủ tuyệt đối chuẩn mực đạo đức \textit{Privacy by Design} (Thiết kế đặt quyền riêng tư làm trọng tâm), hệ thống loại bỏ hoàn toàn các mô hình học sâu (Deep Learning) cồng kềnh trong việc phân tích nội dung. Thay vào đó, phần mở rộng trình duyệt (Chrome Extension) ứng dụng phương pháp "Shift-Left Local NLP Heuristics". Trình duyệt sử dụng một từ điển tĩnh được cấu hình cứng (hardcoded static dictionary) gồm 150 từ khóa mang độc tính cao nhất, được trích xuất ngoại tuyến từ tập dữ liệu Jigsaw Toxic Comment. Tiện ích mở rộng quét toàn bộ cây DOM nội hạt thông qua Biểu thức chính quy (Regular Expressions - Regex) với độ phức tạp thời gian $\mathcal{O}(1)$, từ đó vi tính hóa một Điểm số Độc hại (Toxicity Score, $\Gamma(t)$) có miền giới hạn từ 0.0 đến 1.0 trực tiếp tại thiết bị. Kiến trúc này đảm bảo tuyệt đối không có bất kỳ chuỗi văn bản tự nhiên nào của người dùng (user text) bị truyền tải về máy chủ.

\subsection{Trách nhiệm quản lý: Bộ điều khiển PID (Algorithmic Circuit Breaker)}

Trách nhiệm quản lý phải được tự động hóa. Kiến trúc v3.0 sử dụng nền tảng của lý thuyết điều khiển cổ điển để xây dựng một mạch ngắt thuật toán (algorithmic circuit breaker).

Được truyền cảm hứng từ ứng dụng PID vào kỹ thuật phần mềm của Hellerstein et al. (2004) \cite{hellerstein2004feedback}, bộ điều khiển xử lý biến trạng thái đầu vào:

\begin{equation}
    u(t) = K_p e(t) + K_i \int_{0}^{t} e(\tau)d\tau + K_d \frac{de(t)}{dt}
    \label{eq:pid}
\end{equation}

Khi thao túng tâm lý vọt qua ngưỡng an toàn, bộ điều khiển sinh ra lực ma sát kỹ thuật số (digital friction). Điều này chứng minh tuyệt đối rằng đội ngũ kỹ sư có thể tự tay xây dựng các cơ chế quản lý tự động trực tiếp vào lõi phần mềm.

\subsubsection{Công thức Chỉ số Rủi ro Tổng hợp (Risk Index)}

Tín hiệu kích hoạt mạch ngắt thuật toán (đại lượng $e(t)$ trong khối PID) được định lượng bởi Chỉ số Rủi ro tổng hợp ($I_{risk}(t)$). Chỉ số này liên kết đa chiều các trạng thái tương tác thành một đại lượng phản hồi tuyến tính duy nhất. Hàm rủi ro được thiết lập chính xác bao gồm 3 thành phần cấu thành:

\begin{equation}
    I_{risk}(t) = w_1 \cdot A(t) + w_2 \cdot \Gamma(t) + w_3 \cdot S(t)
    \label{eq:risk_index}
\end{equation}

Trong đó, các biến tham số được định nghĩa như sau:
\begin{itemize}
    \item $A(t)$: Trạng thái Gây nghiện (Addiction) biểu diễn dưới dạng vận tốc cuộn cơ học (Velocity), chịu trọng số cao nhất $w_1 = 0.6$.
    \item $\Gamma(t)$: Điểm số Độc hại (Toxicity) quét cục bộ, đánh giá mức phơi nhiễm theo thời gian thực với trọng số $w_2 = 0.3$.
    \item $S(t)$: Thời lượng Phiên hoạt động (Session Duration) trích xuất sự hao hụt năng lượng sinh lý, vận hành với trọng số $w_3 = 0.1$.
\end{itemize}

\subsubsection{Cơ chế thực thi: Ngụy biện Thiết kế của "Khóa Cứng" (The Fallacy of Hard Blocks)}

Dưới góc độ Kỹ thuật Hệ thống và Thiết kế Giao diện Người-Máy (HCI), việc áp dụng cơ chế "Khóa Cứng" (Hard Blocks) - ví dụ: hiển thị một màn hình cảnh báo cưỡng chế toàn bộ trình duyệt hay chặn hoàn toàn một trang web khi phát hiện sự lạm dụng - là một sai lầm cơ bản mang tính kiến trúc. 

Về mặt tâm lý học hành vi, các biện pháp can thiệp thô bạo này lập tức kích hoạt \textit{Tâm lý Phản kháng} (Psychological Reactance). Khi cảm quan về quyền tự chủ (autonomy) bị đe dọa trực diện, xung đột nhận thức nội tại buộc người xem phản ứng lại theo cách tiêu cực đối với chính hệ thống bảo vệ (tìm cách tắt extension, gỡ cài đặt phần mềm). Khóa cứng là biểu hiện của một triết lý thiết kế hệ thống lười biếng, thiếu khả năng tinh chỉnh độ phân giải của tác động.

Thay vào đó, đầu ra của khối PID Controller trong kiến trúc v3.0 kích hoạt chiến lược \textbf{Ma sát Giao diện Tinh tế} (Mindful Friction), vận hành theo nguyên lý tước đoạt thụ động:

\begin{itemize}
    \item \textbf{Cấp độ 1 - Triệt tiêu thị giác (Desaturation):} Khi $\text{Risk\_Index} > 0.6$, hệ thống tiêm mã CSS nội suy triệt tiêu độ bão hòa màu của toàn bộ bề mặt hiển thị (\textit{grayscale}) theo tỷ lệ tuyến tính với chỉ số rủi ro. Hành động này lặng lẽ thu hồi "dopamine thị giác" do các kỹ sư nền tảng thiết kế sẵn bằng các màu sắc sặc sỡ nhắm vào não bộ vô thức, cắt giảm hàm lượng phần thưởng thu được ($R_t$).
    \item \textbf{Cấp độ 2 - Điều biến Động lượng (Scroll Throttling):} Khi $\text{Risk\_Index} > 0.8$, hệ thống chèn một vi kẽ hở thời gian (\textit{100ms - 300ms delay}) vào bộ xử lý luồng thao tác cuộn (\textit{scroll events}). Can thiệp này chặn đứng độ mượt mà vốn được tối ưu hóa khắt khe (\textit{frictionless}) của luồng thao tác. Nỗ lực để duy trì hành vi lướt tự động (System-1) đột ngột vượt quá mức năng lượng nhận thức khởi điểm, ép buộc não bộ chuyển mạch sang Hệ thống 2 thống trị (System-2) để đánh giá lại hành động cuộn.
\end{itemize}

Tuyệt đối không có bất kỳ dòng cảnh báo dạng văn bản hay nút bấm chặn đứng (pop-up alerts) nào xuất hiện. Ma sát xuất hiện vô hình như lực cản vật lý, trả lại hoàn toàn quyền tự do ý chí nhưng gia tăng tỷ lệ chi phí-lợi ích lên não bộ khi tiêu thụ nội dung.
